\documentclass[11pt]{article}

\usepackage[T1]{fontenc}
\usepackage{amsmath,amssymb,amsthm}
\usepackage[margin=1in]{geometry}
\usepackage[hidelinks]{hyperref}

\newtheorem{theorem}{Theorem}
\newtheorem{lemma}[theorem]{Lemma}
\newtheorem{remark}[theorem]{Remark}

\newcommand{\F}{\mathbb F}
\newcommand{\Tr}{\operatorname{Tr}}
\newcommand{\adj}{\operatorname{adj}}
\newcommand{\supp}{\operatorname{supp}}

\title{Exact Algebraic Degree of Complete Raw Finite-Field Inversion}
\author{Xiaojian Zeng\\
PBC School of Finance, Tsinghua University, 100190\\
\texttt{xiaojianzeng@foxmail.com}}
\date{}
\hypersetup{
  pdftitle={Exact Algebraic Degree of Complete Raw Finite-Field Inversion},
  pdfauthor={Xiaojian Zeng},
  pdfsubject={A uniform exact algebraic-degree theorem for complete raw finite-field inversion},
  pdfkeywords={Boolean functions, algebraic normal form, finite fields, inversion, algebraic degree}
}

\begin{document}
\maketitle

\begin{abstract}
Let $p\in\F_2[t]$ be any irreducible polynomial of degree $n\geq 3$,
let $J_n$ denote inversion in $\F_2[t]/(p)$ in the associated polynomial
basis (with $0^{-1}=0$), and regard the entries of
$P\in\F_2^{n\times n}$ and $u\in\F_2^n$ as independent Boolean variables.
For the complete raw full-domain map
\[
        f_{\rm raw}(P,u)=\adj(P)J_n(Pu),
\]
including singular $P$, we prove
$\deg f_{\rm raw}=3(n-1)$.  The lower bound is witnessed by an explicit
top-degree ANF coefficient.  Its value is an off-diagonal entry of the
inverse trace Gram matrix, while its matrix-variable support is a disjoint
edge together with an alternating spanning path.  This closes the
algebraic-degree half of a two-part question posed by Zhang and Zhang.  It
does not prove the conjectured value of the joint ANF leap, so the original
two-part problem remains unresolved.
\end{abstract}

\section{Context and statement}

Algebraic normal form (ANF) is the unique multilinear polynomial
representation of a Boolean function; algebraic degree is the largest
cardinality of a monomial support with nonzero coefficient.  We use these
standard notions also coordinatewise for vectorial Boolean maps; see
Carlet's Chapter~2 \cite{Carlet2021}.  Standard background on finite fields,
trace, and dual bases may be found in Chapter~2, Section~3 of
\cite{LidlNiederreiter1997}.

Definition~11 of Zhang and Zhang \cite{ZhangZhang2026} introduced the
complete raw polynomial extension above on the full Boolean matrix space.
Their Proposition~12 and Theorem~17 prove the bounds
\[
 \deg f_{\rm raw}\leq 3(n-1),\qquad
 L_{\rm ANF}^{\rm joint}(f_{\rm raw})\geq n,
\]
and Section~4.2 reports exhaustive calculations attaining both bounds for
$n=3,4$.  Their definition is a polynomial extension to every matrix $P$,
not merely the restriction to $P\in\operatorname{GL}_n(\F_2)$.  The contribution here is a
uniform nonzero coefficient that turns the first bound into an equality for
every permitted polynomial reference basis.  No matching upper bound for
the joint ANF leap is asserted.

The witness is obtained by coefficient extraction through two matching
expansions.  Nonzero coefficient extraction over finite fields also appears,
in a different setting, in the permanent-rank argument of Zeng, Liu, and
Ratnavelu \cite{202608.1194}, whose Section~3 combines it with the
Combinatorial Nullstellensatz; that work is cited here
as a related coefficient-certificate method, not as prior work on Boolean
inversion.  The comparison with prior results on the present problem is
limited to the very recent source \cite{ZhangZhang2026}; no global priority
claim is made.

Fix $n\geq3$ and an irreducible $p\in\F_2[t]$ of degree $n$.  Write
\[
 F=\F_2[t]/(p),\qquad \alpha=t\bmod p,
 \qquad E=(e_0,\ldots,e_{n-1})=(1,\alpha,\ldots,\alpha^{n-1}).
\]
Coordinates in $E$ identify $F$ with $\F_2^n$.  Let
\[
 J_n([x]_E)=[x^{-1}]_E,\qquad 0^{-1}=0,
\]
and define $f_{\rm raw}:\F_2^{n\times n}\times\F_2^n\to\F_2^n$ by
\[
 f_{\rm raw}(P,u)=\adj(P)J_n(Pu).
\]
All $n^2+n$ entries of $(P,u)$ are independent Boolean variables.  Thus
the statement concerns the full domain, including singular $P$.

Put
\[
 G_{ij}=\Tr_{F/\F_2}(e_i e_j),\qquad B=G^{-1}.
\]
The trace pairing is nondegenerate, so $G$ and $B$ are nonsingular.

\begin{theorem}\label{thm:main}
For every $n\geq3$ and every irreducible $p\in\F_2[t]$ of degree $n$,
the complete raw full-domain map satisfies
\[
                         \deg f_{\rm raw}=3(n-1).
\]
More precisely, choose distinct indices $a,r$ with $B_{ar}=1$, and choose
distinct column indices $s,k$.  List the columns other than $s$ as
\[
 c_0=k,c_1,\ldots,c_{n-2}
\]
and the rows other than $r$ as
\[
 \rho_0,\ldots,\rho_{n-2}=a.
\]
For
\[
 Q=\{(\rho_i,c_i):0\leq i\leq n-2\}
   \cup\{(r,s)\}
   \cup\{(\rho_i,c_{i+1}):0\leq i<n-2\},                 \tag{1}
\]
the coefficient of
\[
 \left(\prod_{(i,j)\in Q}p_{ij}\right)
 \left(\prod_{j\ne s}u_j\right)                         \tag{2}
\]
in output coordinate $k$ is $B_{ar}=1$.
\end{theorem}

The existence of the off-diagonal entry used in the theorem is proved
below; its location may depend on $p$.

\section{The top layer of coordinate inversion}

Let $(\lambda_0,\ldots,\lambda_{n-1})$ be trace-dual to $E$:
\[
                   \Tr(\lambda_r e_i)=\delta_{ri}.
\]
Since $G$ is symmetric, the coordinate column of $\lambda_r$ is the
$r$th column of $B$:
\[
                         \lambda_r=\sum_a B_{ar}e_a.       \tag{3}
\]

\begin{lemma}\label{lem:top}
The homogeneous degree-$(n-1)$ part of the ANF of $J_n$ is
\[
 (J_n)_{[n-1]}(x)
   =\sum_{r=0}^{n-1}[\lambda_r]_E\prod_{i\ne r}x_i.       \tag{4}
\]
There is no degree-$n$ term.
\end{lemma}

\begin{proof}
For a Boolean map $H:\F_2^n\to\F_2^m$, Mobius inversion gives the
coefficient vector of $\prod_{i\in A}x_i$ as
\[
                         \sum_{C\subseteq A}H(1_C),        \tag{5}
\]
where the sum is in $\F_2^m$.  Hence the field element corresponding to
the coefficient of $\prod_{i\ne r}x_i$ in $J_n$ is
\[
 \sum_{x\in H_r}x^{-1},\qquad
 H_r=\operatorname{span}_{\F_2}\{e_i:i\ne r\}
    =\ker\bigl(x\mapsto\Tr(\lambda_r x)\bigr).            \tag{6}
\]

Let $H_0=\ker\Tr$.  The monic polynomials
\[
 \prod_{h\in H_0}(z+h)
 \quad\hbox{and}\quad
 \Tr(z)=z+z^2+z^4+\cdots+z^{2^{n-1}}
\]
have the same degree and the same $2^{n-1}$ roots, so they are equal.
After division by $z$, comparison of the constant and linear coefficients
gives
\[
                 \sum_{h\in H_0\setminus\{0\}}h^{-1}=1.  \tag{7}
\]
Indeed, the constant coefficient is the product of the nonzero elements of
$H_0$, and the linear coefficient is that product times the sum in (7);
both coefficients in $\Tr(z)/z$ are $1$.

Equation (6) gives $H_r=\lambda_r^{-1}H_0$.  Scaling (7) therefore shows
that (6) is $\lambda_r$, proving all coefficients in (4).  Finally, the
degree-$n$ coefficient corresponds to
\[
  \sum_{x\in F}x^{-1}=\sum_{x\in F}x=0,
\]
because inversion permutes $F\setminus\{0\}$ and the sum of all elements of
$F$ is zero for $n\geq2$.  This proves the lemma.
\end{proof}

\section{An off-diagonal trace-dual coefficient}

\begin{lemma}\label{lem:offdiag}
For a polynomial basis of length $n\geq3$, the inverse trace Gram matrix
$B$ has a nonzero off-diagonal entry.
\end{lemma}

\begin{proof}
If $B$ had no nonzero off-diagonal entry, its nonsingularity over $\F_2$
would force $B=I$, and hence $G=I$.  Let $M$ be the matrix of
multiplication by $\alpha$ in the basis $E$.  Invariance of the trace
pairing under moving the factor $\alpha$ gives
\[
                              M^T G=GM.                   \tag{8}
\]
Thus $G=I$ would make $M$ symmetric.  On the other hand,
$\alpha e_0=e_1$ gives $M_{1,0}=1$, while $\alpha e_1=e_2$ (available
because $n\geq3$) gives $M_{0,1}=0$.  This contradiction proves the claim.
\end{proof}

This argument applies separately to the polynomial basis determined by
every irreducible $p$; it does not assume that the location of the nonzero
off-diagonal entry is basis-invariant.

\section{The alternating-path coefficient}

\begin{proof}[Proof of Theorem~\ref{thm:main}]
Choose $a\ne r$ with $B_{ar}=1$ using Lemma~\ref{lem:offdiag}, and form
$Q$ as in (1).  The target (2) has $2(n-1)$ distinct matrix variables and
$n-1$ distinct vector variables, hence degree $3(n-1)$.

By Lemma~\ref{lem:top}, the top-degree contribution to output coordinate
$k$ is
\[
 \sum_{a',r'} B_{a'r'}\,\adj(P)_{k a'}
                  \prod_{i\ne r'}(Pu)_i.                 \tag{9}
\]
In a term of $\prod_{i\ne r'}(Pu)_i$, the coefficient of a squarefree
$u$-monomial of degree $n-1$ comes from assigning distinct $u$-columns to
the $n-1$ rows other than $r'$.  Thus a term producing
$\prod_{j\ne s}u_j$ selects a perfect matching from the rows other than
$r'$ to the columns other than $s$.  Likewise, the determinant expansion
of $\adj(P)_{k a'}$ selects a perfect matching from the rows other than
$a'$ to the columns other than $k$.  Signs disappear over $\F_2$.

To produce all $2(n-1)$ distinct matrix variables in (2), these two
matchings must be edge-disjoint and have union $Q$.  Regard $Q$ as a
bipartite row--column graph.  It is the disjoint union of the edge $s-r$
and the path
\[
 k=c_0-\rho_0-c_1-\rho_1-\cdots-c_{n-2}-\rho_{n-2}=a.    \tag{10}
\]
The matching arising from $J_n(Pu)$ cannot use column $s$, so the isolated
edge $s-r$ must lie in the adjugate matching.  The adjugate matching cannot
use column $k$, so the first edge of (10) must lie in the inversion
matching.  The matching constraints then force alternation along the whole
path.  Consequently the decomposition is unique.  The inversion matching
is
\[
                    \{(\rho_i,c_i):0\leq i\leq n-2\},
\]
which omits row $r$, and the adjugate matching is
\[
       \{(r,s)\}\cup\{(\rho_i,c_{i+1}):0\leq i<n-2\},
\]
which omits row $a$.  The coefficient of (2) in (9) is therefore exactly
$B_{ar}=1$.

No lower ANF layer of $J_n$ can contribute to (2): a degree-$d$ ANF
monomial has degree at most $2d$ after the bilinear substitution $x=Pu$,
and multiplication by an adjugate entry adds at most $n-1$.  When
$d<n-1$, this is strictly less than $3(n-1)$.  Hence (2) is genuinely a
nonzero ANF monomial of degree $3(n-1)$.

Conversely, every adjugate entry has degree at most $n-1$, and $J_n$ has
degree $n-1$ (also immediate from Lemma~\ref{lem:top}).  Substitution of
the bilinear coordinates of $Pu$ at most doubles the latter degree, so
\[
                  \deg f_{\rm raw}\leq(n-1)+2(n-1).
\]
The displayed coefficient supplies the reverse inequality and completes
the proof.
\end{proof}

\begin{remark}[Dependence on the reference polynomial]
If a second basis is expressed in the first by $T\in\operatorname{GL}_n
(\F_2)$, then
\[
                  G'=T^TGT,\qquad B'=T^{-1}BT^{-T}.
\]
Thus the particular witness location $(a,r)$ need not survive a basis
change.  Uniformity means instead that Lemma~\ref{lem:offdiag} applies in
each allowed polynomial basis, after which the witness is reselected.
\end{remark}

\begin{remark}[Unresolved part]
The source \cite{ZhangZhang2026} defines the joint ANF leap using an
optimal ordering of the joint monomial support.  The coefficient family
above proves no ordering that introduces at most $n$ new variables at each
step.  Hence it does not establish
$L_{\rm ANF}^{\rm joint}(f_{\rm raw})\leq n$; the equality for general
$n$ and the original conjunction remain unresolved.
\end{remark}

\section{Exact finite checks and assistance disclosure}

A deterministic Python verifier computes the coefficient in (2) directly
from the Boolean Mobius formula (5).  It enumerates every monic irreducible
binary polynomial of degrees $3$, $4$, and $5$: respectively $2$, $3$, and
$6$ polynomials.  All eleven computed coefficients are $1$.  These are
exact finite checks of the displayed certificate, not evidence replacing
the quantified proof.  The same proof-versus-verifier separation is made
explicit in the related coefficient-certificate workflow of
\cite{202608.1194}.

An OpenAI Codex language model assisted with organizing the proof, checking
its scope against the supplied source and referee report, drafting the
manuscript, and preparing the exact verification workflow.  The author is
responsible for the mathematical claims.  The uniform theorem rests on the
symbolic argument above; the finite computations cover only the stated
degrees.

\begingroup
\raggedright
\bibliographystyle{plain}
\bibliography{references}
\endgroup

\end{document}
